\subsection{可解李代数的定义}

定义 1。若李代数 g 的第 \(\mathcal{D}^{k}\mathfrak{g}\) 个导出代数 \(k\) 在 k 充分大时为零，则称 g 是可解的。

幂零李代数是可解的。

命题 1. 可解李代数的子代数和商代数都是可解的。可解代数由可解代数扩张得到的每个扩张都是可解的。有限个可解代数的直积也是可解的。

设 g 是李代数，\(g'\) 是子代数，h 是 g 的理想，t = g/h，\(\phi\) 是从 g 到 t 的典范映射。若 g 可解，则对某个整数 k 有 \(D^{k}g = \{0\}\)，从而 \(D^{k}g' \subset D^{k}g = \{0\}\)，\(D^{k}t = \phi(\mathcal{D}^{k}g) = \{0\}\)，所以 \(g'\) 和 t 都可解。若 h 和 t 可解，则存在整数 s、t，使得

\[
\mathcal {D} ^ {s} \mathfrak {h} = \mathcal {D} ^ {t} \mathfrak {k} = \{0 \};
\]

则 \(\mathcal{D}^t\mathfrak{g}\subset \mathfrak{h}\)，从而 \(\mathcal{D}^{s + t}\mathfrak{g} = \mathcal{D}^{s}(\mathcal{D}^{t}\mathfrak{g})\subset \mathcal{D}^{s}\mathfrak{h} = \{0\}\)，且 \(\mathfrak{g}\) 可解。最后一个断言由第二个断言对因子个数作归纳得到。

命题 2. 设 g 是李代数。以下条件等价：

\begin{enumerate}
\def\labelenumi{(\alph{enumi})}
\item
  g 是可解的；
\item
  存在 \(\mathfrak{g} = \mathfrak{g}_0 \supset \mathfrak{g}_1 \supset \dots \supset \mathfrak{g}_n = \{0\}\) 的理想递降序列 \(\mathfrak{g}\)，使得代数 \(\mathfrak{g}_i / \mathfrak{g}_{i+1}\) 是交换的（\(i = 0, 1, \ldots, n-1\)）；
\item
  存在 \(\mathfrak{g} = \mathfrak{g}_0' \supset \mathfrak{g}_1' \supset \dots \supset \mathfrak{g}_p' = \{0\}\) 的子代数递降序列 \(\mathfrak{g}\)，使得 \(\mathfrak{g}_{i+1}'\) 是 \(\mathfrak{g}_i'\) 的理想且 \(\mathfrak{g}_i'/\mathfrak{g}_{i+1}'\) 是交换的（\(i = 0, 1, \ldots, p - 1\)）；
\item
  存在 \(\mathfrak{g} = \mathfrak{g}_0'' \supset \mathfrak{g}_1'' \supset \dots \supset \mathfrak{g}_q'' = \{0\}\) 的子代数递降序列 \(\mathfrak{g}\)，使得 \(\mathfrak{g}_{i+1}''\) 是 \(\mathfrak{g}_i''\) 的余维 1 理想（\(i = 0, 1, \ldots, q - 1\)）。
\item
  \(\Rightarrow\) (b)：只需考虑 g 的导出理想序列。
\item
  \(\Rightarrow\) (c)：显然。
\item
  \(\Rightarrow\) (d)：假设条件 (c) 成立；\(\mathfrak{g}_i'\) 中每个包含 \(\mathfrak{g}_{i+1}'\) 的向量子空间都是 \(\mathfrak{g}_i'\) 的理想，故 (d) 立即成立。
\item
  \(\Rightarrow\) (a)：这立即由可解代数由可解代数扩张所得的扩张仍可解这一事实得到。
\end{enumerate}

